Behovet av att optimera lagring och överföring av digitala bilder har lett till utvecklingen av en mängd olika algoritmer vilka generellt delas in i förlustfria och förstörande metoder. Förlustfritt format som PNG (Portable Network Graphics), bevarar bildens exakta pixeldata genom att utnyttja statistiska överflöd vilket är avgörande i sammanhang där ingen information får gå förlorad [4]. Å andra sidan använder de format såsom den klassiska JPG-standarden och det modernare AVIF tekniker för att ta bort visuell information som det mänskliga ögat har svårt att uppfatta, vilket möjliggör en dramatisk minskning av filstorleken[7]. AVIF (AV1 Image File Format) är särskilt relevant för denna studie eftersom det bygger på modern videokodningsteknik och erbjuder stöd för både högt bitdjup och extra kanaler såsom alfakanalen för transparens (RGBA), vilket är nödvändigt för att hantera de latenta representationernas fyra kanaler [3]. Utvecklingen av nyare format som JPEG XL (JXL) fortsätter att pressa gränserna för hur effektivt komplex bildinformation kan representeras utan att den upplevda kvaliteten försämras.